\begin{afterword}
\summaryhead

The essay defines rationality in two parts. Epistemic rationality is systematically improving the accuracy of your beliefs; instrumental rationality is systematically achieving your values, which the author calls ``winning.'' Truth here does not mean certainty, and winning does not mean beating other people.

The author argues that the word ``rational'' is useful for the same reason ``true'' is: it lets us state general facts about good ways of thinking. Psychologists judge reasoning against probability theory and decision theory; someone who rates ``Bill is an accountant who plays jazz'' above ``Bill plays jazz'' breaks a law of probability. But these formal standards cannot be computed in real life, and their meaning is disputed in some puzzles, so there is an art of thinking well from inside a human mind.

The post denies that it is an argument about a word. If you find yourself saying that something is rational but false, or rational but wrong, use more specific language. The essay warns against overusing the word.

\responsehead

The essay denies that it is arguing about the meaning of a word. It spends most of its length defining one, and at two points it uses the definition in place of an argument.

The definitions are the standard pair from philosophy, presented without their known difficulty: what should you do when a false belief would serve your values? The essay never asks. Its rule for belief is sound as far as it goes: one cannot coherently hold, at the same moment, that X is rational to believe and that the truth is Y. But it says nothing about the ordinary backward-looking case, a belief that was rational on misleading evidence and turned out false. For action it goes further, and declares any gap between the rational and the right ``almost certainly'' a confusion of words, with no reason given.

The Newcomb example takes a side without saying so. ``I'd rather have a million dollars'' presents one side of an open dispute as plain common sense. The essay's own cited source says people split almost evenly on the problem, and the philosophers surveyed in 2020 leaned the other way. A few paragraphs earlier the essay had itself listed Newcomblike problems among the cases where the mathematics is in dispute.

By the end the author describes using ``rationality'' ``normatively, to pick out desirable patterns of thought.'' That fits the essay's answer to its opening objection, that the word is needed to state general facts about good thinking. But in any single case the word then adds nothing to ``good,'' and the vocabulary built on it (Bayesian, winning, the Way) expresses approval without adding content.

There is little practical content. The essay concedes that no one can calculate and obey the mathematics it names as the ``gold standard,'' and never says what guides a person in its place. What it offers instead is a list of aims (learn your flaws, overcome biases, get into good emotional shape) that ends ``et cetera, et cetera.''

In short: an essay that says it is not arguing about a word, spends most of its length defining one, takes a side on Newcomb's problem without saying the side is contested, and ends its description of the practical art with ``et cetera, et cetera.''
\end{afterword}
